\section{Introduction}

La planification de maintenance ferroviaire représente un défi opérationnel majeur pour les entreprises de transport, notamment la SNCF. Avec des centaines de trains nécessitant des maintenances régulières, des fenêtres temporelles strictes, et des ressources limitées (voies, techniciens, équipements), ce problème présente toutes les caractéristiques d'un problème d'optimisation combinatoire complexe.

Ce projet s'intéresse à la modélisation de ce problème comme un \textbf{Resource-Constrained Project Scheduling Problem (RCPSP)} avec des contraintes spécifiques au domaine ferroviaire : fenêtres temporelles strictes (possiblement multiples par train), temps de setup variables selon le contexte, et ressources multiples avec capacités simultanées.

\subsection{Problématique}

Le problème consiste à planifier les opérations de maintenance sur un horizon de plusieurs jours en respectant :
\begin{itemize}
	\item Les fenêtres temporelles de présence des trains au centre de maintenance
	\item Les contraintes de précédence entre opérations avec temps de setup
	\item Les capacités limitées des ressources (voies exclusives, techniciens par type, équipements)
	\item Les contraintes de travail des techniciens (temps min/max, pauses)
	\item Les compatibilités entre voies et équipements
	\item La gestion dynamique des deadlines et des maintenances optionnelles
\end{itemize}

L'objectif est de minimiser le retard pondéré des maintenances par rapport à leurs dates butoirs, calculées en fonction de leur périodicité et de la dernière fois où elles ont été effectuées.

\subsection{Objectifs}

Les objectifs de ce projet sont :
\begin{enumerate}
	\item Modéliser le problème de planification de maintenance ferroviaire comme un RCPSP avec contraintes temporelles
	\item Développer deux modèles MiniZinc distincts : l-v1 avec vérification ponctuelle des ressources, et e-v1 avec modélisation par blocs de travail
	\item Comparer les performances des deux modèles sur trois solveurs (Chuffed, Gecode, OR-Tools) sur différentes instances
	\item Identifier les limites de résolvabilité selon la complexité des instances
	\item Analyser l'impact des choix de modélisation (vérification ponctuelle vs blocs de travail) sur la performance
\end{enumerate}

\subsection{Structure du rapport}

Ce rapport est organisé comme suit : la Section~\ref{sec:probleme} présente la description formelle du problème, la Section~\ref{sec:approche} décrit les deux modèles développés et leurs différences, la Section~\ref{sec:protocole} présente le protocole expérimental, la Section~\ref{sec:resultats} expose les résultats expérimentaux, la Section~\ref{sec:discussion} en discute les implications, et la Section~\ref{sec:conclusion} conclut le travail.
